\chapter{Solid State Theory: Final, June 24, 2020}
\label{q:theory-2020-06-final}
\begin{sourceNote}
The supplied printed page contains Questions 1--12. Its visible weights sum to 85\%; they are retained below. Handwritten responses are excluded from the transcription.
\end{sourceNote}
\begin{enumerate}[label=\textbf{\arabic*.}]
\item (5\%) Describe the main differences between the Hartree theory and the Hartree--Fock theory.
\item (10\%) From the expression of the energy of the normal ground state of the homogeneous electron gas in the Hartree--Fock approximation
\[
E_0^{(\mathrm N)}=N\left(\frac35\frac{\hbar^2k_F^2}{2m}-\frac34\frac{e^2k_F}{\pi}\right),
\]
express the energy of the fully spin polarized ground state and explain your answer.
\item (10\%) Write down the Hartree and the Hartree--Fock equations and explain each term in these equations. Describe how to solve the equations.
\item (5\%) Describe the advantage of the density functional theory over the Hartree--Fock theory.
\item (5\%) Express and explain the Hohenberg--Kohn functional.
\item (5\%) Express and explain the Kohn--Sham equations.
\item (5\%) Describe the connection of the Kohn--Sham equations with the Hohenberg--Kohn functional.
\item (5\%) Given $\varepsilon_{\mathrm{xc}}(n)$, derive an expression for $V_{\mathrm{xc}}^{(\mathrm{LDA})}(\vct r)$, where $\varepsilon_{\mathrm{xc}}(n)$ is the many-electron exchange-correlation energy per electron of a homogeneous interacting electron gas of electron density $n$, and $V_{\mathrm{xc}}^{(\mathrm{LDA})}(\vct r)$ is the exchange-correlation potential in the local density approximation.
\item (5\%) Describe the Bloch sums in the tight-binding method for the band structure calculation of solids.
\item (15\%) The electronic configuration of the silicon atom is $1s^2 2s^2 2p^6 3s^2 3p^2$.
\begin{enumerate}[label=(\alph*)]
\item (5\%) If five orbitals for each atom are chosen to construct the Bloch sums in the tight-binding method of band structure calculation for electrical and chemical properties of crystalline silicon, what is your suggestion on the five orbitals?
\item (5\%) From (a), how many bands will you have?
\item (5\%) From (b), how many bands will be fully filled by valence electrons for the silicon crystal in the ground state?
\end{enumerate}
\item (10\%) What is the difficulty in band structure calculation of solids using plane waves as basis functions? How to solve the difficulty?
\item (5\%) Describe the core Bloch sum(s) in the orthogonalized plane waves for calculation of the band structure of crystal lithium hydride (LiH).
\end{enumerate}
